\documentclass[10pt, letterpaper]{article}
% Packages:
\usepackage[
    ignoreheadfoot, % set margins without considering header and footer
    top=2 cm, % seperation between body and page edge from the top
    bottom=2 cm, % seperation between body and page edge from the bottom
    left=2 cm, % seperation between body and page edge from the left
    right=2 cm, % seperation between body and page edge from the right
    footskip=1.0 cm, % seperation between body and footer
    % showframe % for debugging 
]{geometry} % for adjusting page geometry
\usepackage{titlesec} % for customizing section titles
\usepackage{tabularx} % for making tables with fixed width columns
\usepackage{array} % tabularx requires this
\usepackage[dvipsnames]{xcolor} % for coloring text
\definecolor{primaryColor}{RGB}{0, 79, 144} % define primary color
\usepackage{enumitem} % for customizing lists
\usepackage{fontawesome5} % for using icons
\usepackage{amsmath} % for math
\usepackage[
    pdftitle={Matteo Romilio Pasqual's CV},
    pdfauthor={Matteo Romilio Pasqual},
    pdfcreator={LaTeX with RenderCV},
    colorlinks=true,
    urlcolor=primaryColor
]{hyperref} % for links, metadata and bookmarks
\usepackage[pscoord]{eso-pic} % for floating text on the page
\usepackage{calc} % for calculating lengths
\usepackage{bookmark} % for bookmarks
\usepackage{lastpage} % for getting the total number of pages
\usepackage{changepage} % for one column entries (adjustwidth environment)
\usepackage{paracol} % for two and three column entries
\usepackage{ifthen} % for conditional statements
\usepackage{needspace} % for avoiding page brake right after the section title
\usepackage{iftex} % check if engine is pdflatex, xetex or luatex
% Ensure that generate pdf is machine readable/ATS parsable:
\ifPDFTeX
    \input{glyphtounicode}
    \pdfgentounicode=1
    % \usepackage[T1]{fontenc} % this breaks sb2nov
    \usepackage[utf8]{inputenc}
    \usepackage{lmodern}
\fi
% Some settings:
\AtBeginEnvironment{adjustwidth}{\partopsep0pt} % remove space before adjustwidth environment
\pagestyle{empty} % no header or footer
\setcounter{secnumdepth}{0} % no section numbering
\setlength{\parindent}{0pt} % no indentation
\setlength{\topskip}{0pt} % no top skip
\setlength{\columnsep}{0cm} % set column seperation
\makeatletter
\let\ps@customFooterStyle\ps@plain % Copy the plain style to customFooterStyle
\patchcmd{\ps@customFooterStyle}{\thepage}{
    \color{gray}\textit{\small Matteo Romilio Pasqual - Page \thepage{} of \pageref*{LastPage}}
}{}{} % replace number by desired string
\makeatother
\pagestyle{customFooterStyle}
\titleformat{\section}{\needspace{4\baselineskip}\bfseries\large}{}{0pt}{}[\vspace{1pt}\titlerule]
\titlespacing{\section}{-1pt}{0.3 cm}{0.2 cm} % section title spacing
\renewcommand\labelitemi{$\circ$} % custom bullet points
\newenvironment{highlights}{
    \begin{itemize}[
        topsep=0.10 cm,
        parsep=0.10 cm,
        partopsep=0pt,
        itemsep=0pt,
        leftmargin=0.4 cm + 10pt
    ]
}{
    \end{itemize}
} % new environment for highlights
\newenvironment{highlightsforbulletentries}{
    \begin{itemize}[
        topsep=0.10 cm,
        parsep=0.10 cm,
        partopsep=0pt,
        itemsep=0pt,
        leftmargin=10pt
    ]
}{
    \end{itemize}
} % new environment for highlights for bullet entries
\newenvironment{onecolentry}{
    \begin{adjustwidth}{
        0.2 cm + 0.00001 cm
    }{
        0.2 cm + 0.00001 cm
    }
}{
    \end{adjustwidth}
} % new environment for one column entries
\newenvironment{twocolentry}[2][]{
    \onecolentry
    \def\secondColumn{#2}
    \setcolumnwidth{\fill, 4.5 cm}
    \begin{paracol}{2}
}{
    \switchcolumn \raggedleft \secondColumn
    \end{paracol}
    \endonecolentry
} % new environment for two column entries
\newenvironment{header}{
    \setlength{\topsep}{0pt}\par\kern\topsep\centering\linespread{1.5}
}{
    \par\kern\topsep
} % new environment for the header
\newcommand{\placelastupdatedtext}{% \placetextbox{<horizontal pos>}{<vertical pos>}{<stuff>}
  \AddToShipoutPictureFG*{% Add <stuff> to current page foreground
    \put(
        \LenToUnit{\paperwidth-2 cm-0.2 cm+0.05cm},
        \LenToUnit{\paperheight-1.0 cm}
    ){\vtop{{\null}\makebox[0pt][c]{
        \small\color{gray}\textit{Last updated in September 2024}\hspace{\widthof{Last updated in September 2024}}
    }}}%
  }%
}%
% save the original href command in a new command:
\let\hrefWithoutArrow\href
% new command for external links:
\renewcommand{\href}[2]{\hrefWithoutArrow{#1}{\ifthenelse{\equal{#2}{}}{ }{#2 }\raisebox{.15ex}{\footnotesize \faExternalLink*}}}
\begin{document}
\newcommand{\AND}{\unskip\cleaders\copy\ANDbox\hskip\wd\ANDbox\ignorespaces}
\newsavebox\ANDbox
\sbox\ANDbox{}

%--------------------------HEADER----------------------------------------------
\begin{header}
    \textbf{\fontsize{24 pt}{24 pt}\selectfont Matteo Romilio Pasqual}
        
    \vspace{0.3 cm}
        
    \normalsize
    \mbox{{\color{black}\footnotesize\faMapMarker*}\hspace*{0.13cm}Milan (IT)}%
    \kern 0.25 cm%
    \AND%
    \kern 0.25 cm
    \mbox{\hrefWithoutArrow{mailto:matteoromilio.pasqual@mail.polimi.it}{\color{black}{\footnotesize\faEnvelope[regular]}\hspace*{0.13cm}matteoromilio.pasqual@mail.polimi.it}}%
    \kern 0.25 cm
    \AND
    \kern 0.25 cm
    \mbox{\hrefWithoutArrow{https://www.linkedin.com/in/matteo-romilio-pasqual/}{\color{black}{\footnotesize\faLinkedinIn}\hspace*{0.13cm}Matteo Pasqual}}

    "In the middle of difficulty lies opportunity." – Albert Einstein
        
\end{header}
\vspace{0.3 cm - 0.3 cm}

%--------------------------EDUCATION----------------------------------------------
\section{\mbox{{\color{black}\footnotesize\faGraduationCap}\hspace*{0.13cm}Education}}
    
    \begin{twocolentry}{    
    \textit{Sep 2024 – Today}}
    \textbf{Politecnico di Milano}
    M.Sc. in Computer Science and Engineering - Artificial intelligence
    \end{twocolentry}   

    \vspace{0.10 cm}
    \begin{twocolentry}{
    \textit{Sep 2021 – Jul 2024}}
    \textbf{Politecnico di Milano}
    \end{twocolentry}
    \begin{onecolentry}
    B.Sc. in Computer Science and Engineering \hspace{1pt}(\href{https://www.polimi.it/}{polimi.it}) \hspace{76pt}Grade: 109/110
    \end{onecolentry}

    \vspace{0.10 cm}
    \begin{twocolentry}{
    \textit{Sep 2016 – Jul 2021}}
    \textbf{Liceo Scientifico "Leonardo da Vinci"}
    \end{twocolentry}
    \begin{onecolentry}
    Diploma di Liceo Scientifico - Scienze Applicate \hspace{1pt}(\href{https://www.liceogallarate.edu.it/} {liceogallarate.edu.it}) \hspace{12pt}Grade: 96/100
    \end{onecolentry}



%--------------------------RESEARCH----------------------------------------------
\section{\mbox{{\color{black}\footnotesize\faReact}\hspace*{0.13cm}Projects and Research}}

\vspace{0.10 cm}
\begin{twocolentry}{\textit{Ongoing Thesis - 2026}}
\textbf{SHRED-Based Sensor Reconstruction in High-Dimensional Dynamical Systems with Missing and Irregular Measurements}
\end{twocolentry}

\vspace{0.10 cm}
\begin{onecolentry}
\begin{highlights}
    \item Investigating SHRED-based architectures for state reconstruction in high-dimensional dynamical systems under missing, asynchronous, and unreliable sensor measurements
    \item Studying the impact of incomplete temporal coverage on recurrent and attention-based encoders within the SHRED framework developing and benchmarking strategies for handling missing data, including imputation methods and masking mechanisms for sequential models
\end{highlights}
\end{onecolentry}

\vspace{0.20 cm}
\begin{twocolentry}{\textit{Project - 2025}}
    \textbf{Multi-agent System Leonardo Drone Contest 2025}
\end{twocolentry}
\vspace{0.10 cm}
\begin{onecolentry}
    \begin{highlights}         
        \item Participated in a national robotics competition (\href{https://www.leonardo.com/it/innovation-technology/open-innovation/drone-contest}{Leonardo Drone Contest})
        \item Designed a multi-agent system for collaborative navigation and dynamic task allocation in unknown indoor environments, developing decentralized decision-making strategies, inter-agent communication protocols, and computer vision modules for object detection and pose estimation
        \item Strengthened skills in teamwork, robustness engineering, and problem-solving under uncertainty
    \end{highlights}
\end{onecolentry}

\vspace{0.20 cm}
\begin{twocolentry}{\textit{Project - 2025}}
    \textbf{Course Projects}
\end{twocolentry}

\vspace{0.10 cm}
\begin{onecolentry}
    \begin{highlights}
        \item ANNDL course project: multivariate time-series classification using neural networks (\href{https://github.com/Gradient-Gang/ANN-Challenges.git}{GitHub})
        \item ANNDL course project: biomedical image classification using deep learning methods (\href{https://github.com/Gradient-Gang/ANN-Challenge-2.git}{GitHub})
        \item MMMIP course project: image processing for denoising, inpainting, and anomaly detection (\href{https://github.com/matteopasqual02/MMMIP-2025-homeworks.git}{GitHub})
        \item OLA course project: pricing under budget constraints using online learning methods (\href{https://github.com/matteopasqual02/OLA-2025-project.git}{GitHub})
    \end{highlights}
\end{onecolentry}

\vspace{0.20 cm}
\begin{twocolentry}{\textit{Research - 2024}}
    \textbf{Kolmogorov-Arnold Networks (KAN) lecture}
\end{twocolentry}
\vspace{0.10 cm}
\begin{onecolentry}
    \begin{highlights}
    
        \item Delivered a lecture on KANs, a neural architecture inspired by the Kolmogorov–Arnold representation theorem, as part of the Numerical Analysis for Machine Learning course
        \item Presented the theoretical foundations of KANs and their differences from MLP, focusing on spline-based learnable activation functions, training methodology, sparsification, pruning, and symbolic regression techniques, highlighting advantages in interpretability and function approximation
        \item Achieved top evaluation (30/30) and selected for inclusion in future NAML lectures at PoliMi
        \item Developed supporting materials and implementation examples (\href{https://github.com/matteopasqual02/KAN-naml.git}{GitHub})
    \end{highlights}
\end{onecolentry}

\vspace{0.20 cm}
\begin{twocolentry}{\textit{Research - 2024}}
    \textbf{Non-intrusive Monitoring of Daily Habits in Older Adults: \\ A Chaos Game Representation-Based Approach}
\end{twocolentry}
\vspace{0.10 cm}
\begin{onecolentry}
    \begin{highlights}
        \item Co-authored a paper accepted at UCAmI 2025 – 17th International Conference on Ubiquitous Computing and Ambient Intelligence (\href{https://doi.org/10.1007/978-3-032-16992-1_20}{Paper})
        \item Developed a novel framework using Chaos Game Representation (CGR) to transform time-series sensor data into fractal-like images for behavioural analysis, enabling classification of daily routines.
        \item Demonstrated effective anomaly detection and clustering of behavioural patterns on both synthetic and real-world multi-resident datasets
        
    \end{highlights}
\end{onecolentry}

%--------------------------EXPERIENCE----------------------------------------------
\section{\mbox{{\color{black}\footnotesize\faShoePrints}\hspace*{0.13cm}Experience}}

    \begin{twocolentry}{
    \textit{Milano, IT}    
            
    \textit{Sep 2025 – Jul 2026}}
    \textbf{University Multichanche  Tutor}
            
    \textit{Politecnico di Milano}
    \end{twocolentry}
    \vspace{0.10 cm}
    \begin{onecolentry}
    \begin{highlights}
        \item Tutored university students with the sponsorship of the Multichance Team
        \item Provided personalized academic support and study assistance

    \end{highlights}
    \end{onecolentry}

\vspace{0.20 cm}

    \begin{twocolentry}{
    \textit{Milano, IT}    
            
    \textit{Sep 2022 – Jul 2024}}
    \textbf{University Peer-to-peer Tutor}
            
    \textit{Politecnico di Milano}
    \end{twocolentry}
    \vspace{0.10 cm}
    \begin{onecolentry}
    \begin{highlights}
        \item Tutored university students with the sponsorship of Univeristy
        \item Provided personalized academic support and study assistance
    \end{highlights}
    \end{onecolentry}





%--------------------------LANGUAGES----------------------------------------------    
\section{\mbox{{\color{black}\footnotesize\faGlobeAmericas}\hspace*{0.13cm}Languages}}

\begin{highlights}
    \item \textbf{Italian}: native
    \item \textbf{English}: C1
\end{highlights}

%--------------------------SKILLS----------------------------------------------
\section{\mbox{{\color{black}\footnotesize\faWrench}\hspace*{0.13cm}Skills}}

    \begin{samepage}
    \begin{highlights}
        \item \textbf{Programming \& Software Dev.}: Python, C/C++, Java
        \item \textbf{Databases Technologies}: SQL, MongoDB, Spark, Neo4j, Zilliz-Milvus
        \item \textbf{OS \& Middelware}: ROS, ROS2, Linux, Docker and VMs
        \item \textbf{Other Technologies}: Matlab, VHDL/Vivado
    \end{highlights}
    \end{samepage}

%--------------------------AWARDS----------------------------------------------  
\section{\mbox{{\color{black}\footnotesize\faTrophy}\hspace*{0.13cm}Awards}}
        \begin{twocolentry}{2023}
            \textbf{Best Freshmen Award}: top 5\% freshmen of PoliMi
        \end{twocolentry}

%--------------------------Volounteering---------------------------------------------- 
    \section{\mbox{{\color{black}\footnotesize\faSeedling}\hspace*{0.13cm}Volounteering}}
            
        \begin{twocolentry}{Set 2016 - Aug 2021}
        \textbf{Volounteer at Oratorio di Albizzate}\\
        Albizzate
        \end{twocolentry}

\end{document}





